
\documentclass{article}
\usepackage{colortbl}
\usepackage{makecell}
\usepackage{multirow}
\usepackage{supertabular}

\begin{document}

\newcounter{utterance}

\centering \large Interaction Transcript for game `dond', experiment `semi', episode 14 with Qwen3.6{-}35B{-}A3B{-}FP8{-}without{-}reasoning.
\vspace{24pt}

{ \footnotesize  \setcounter{utterance}{1}
\setlength{\tabcolsep}{0pt}
\begin{supertabular}{c@{$\;$}|p{.15\linewidth}@{}p{.15\linewidth}p{.15\linewidth}p{.15\linewidth}p{.15\linewidth}p{.15\linewidth}}
    \# & $\;$A & \multicolumn{4}{c}{Game Master} & $\;\:$B\\
    \hline

    \theutterance \stepcounter{utterance}  
    & & \multicolumn{4}{p{0.6\linewidth}}{
        \cellcolor[rgb]{0.9,0.9,0.9}{
            \makecell[{{p{\linewidth}}}]{
                \texttt{\tiny{[P1$\langle$GM]}}
                \texttt{Mängite koostöö läbirääkimiste mängu, milles peate saavutama teise mängijaga kokkulepet selle kohta, kuidas objektide komplekti teie vahel jagada.} \\
\\ 
\texttt{Reeglid:} \\
\texttt{(a) Teie ja teine mängija saate objektide komplekti. Igaühele antakse ka salajane väärtuse funktsioon, mis näitab, kui palju hindasite iga objekti tüüpi.} \\
\texttt{(b) Vahetate teise mängijaga sõnumeid, et kokku leppida selles, kes saab milliseid objekte. Igaüks võib saata kuni 5 sõnumit või lõpetada mängu varajalt, tehes igal ajal salajase ettepaneku.} \\
\texttt{(c) Teil palutakse esitada salajane ettepanek, mis näitab objekte, mida soovite, vormindatuna ruuduliste sulgudega järgmiselt: "{[}Ettepanek: <arv> <objekti nimi>, <arv> <objekti nimi>, <...>{]}"} \\
\texttt{(d) Kui teie ettepanekud on täiendavad, st ettepaanekute täitmiseks on piisavalt objekte, saab iga mängija skoori vastavalt tema saanud objektide väärtuste summale. Muul juhul saavad mõlemad nullpunkte.} \\
\texttt{(e) Teie eesmärk on maksimeerida oma skoori ja teise mängija skoori summa.} \\
\\ 
\texttt{Alustame.} \\
\\ 
\texttt{Saadaolevate objektide komplekt on:} \\
\\ 
\texttt{3 poldid, 1 tellis, 1 lamp} \\
\\ 
\texttt{Teie salajane väärtuse funktsioon on:} \\
\\ 
\texttt{polt: 1, tellis: 4, lamp: 3} \\
\\ 
\texttt{OLULINE: Teie sõnumid, välja arvatud salajane ettepanek, edastatakse otse teisele mängijale, nii et ärge lisage mingit vastust reeglitele ega teksti, mis teie sõnumi kuulutaks. Salajase ettepaneku tegemiseks kasutage määratletud vormingut. Ärge kasutage ruudulisi sulge, kui suhtlesite teise mängijaga, või seda interpreteeritakse teie salajase ettepanekuna.} \\
            }
        }
    }
    & \\ \\

    \theutterance \stepcounter{utterance}  
    & \multicolumn{4}{p{0.6\linewidth}}{
        \cellcolor[rgb]{0.9,0.9,0.9}{
            \makecell[{{p{\linewidth}}}]{
                \texttt{\tiny{[P1$\rangle$GM]}}
                \texttt{Tere! Mul on hea meel koostööd teha.} \\
\\ 
\texttt{Minu hinnangute põhjal on mulle kõige väärtuslikum tellis (4 punkti), seejärel lamp (3 punkti) ja poltid (1 punkt).} \\
\\ 
\texttt{Kuna meie eesmärk on maksimeerida mõlema mängija skooride summat, proovime leida lahenduse, kus mõlemad saavad oma hinnangute põhjal kõige väärtuslikumad objektid.} \\
\\ 
\texttt{Kas teil on juba mõte, kuidas jagame? Kui jah, palun jagage oma ettepanek või hinnangud. Kui ei, siis ma pakun välja järgmise jaotuse, mis tundub mulle mõlemale kasulik:} \\
\\ 
\texttt{{-} **Mulle:** tellis ja lamp (kokku 7 punkti minu jaoks)} \\
\texttt{{-} **Teile:** 3 poldi (3 punkti teie jaoks, eeldades, et teie hinnangud on sarnased või erinevad, aga poldid on tavaliselt odavamad)} \\
\\ 
\texttt{Kas see jaotus sobib teile? Või soovite teha oma ettepaneku?} \\
            }
        }
    }
    & & \\ \\

    \theutterance \stepcounter{utterance}  
    & & & \multicolumn{2}{p{0.3\linewidth}}{
        \cellcolor[rgb]{0.9,0.9,0.9}{
            \makecell[{{p{\linewidth}}}]{
                \texttt{\tiny{[GM$|$GM]}}
                \texttt{continue} \\
            }
        }
    }
    & & \\ \\

    \theutterance \stepcounter{utterance}  
    & & & \multicolumn{2}{p{0.3\linewidth}}{
        \cellcolor[rgb]{0.9,0.9,0.9}{
            \makecell[{{p{\linewidth}}}]{
                \texttt{\tiny{[GM$|$GM]}}
                \texttt{Tere! Mul on hea meel koostööd teha.} \\
\\ 
\texttt{Minu hinnangute põhjal on mulle kõige väärtuslikum tellis (4 punkti), seejärel lamp (3 punkti) ja poltid (1 punkt).} \\
\\ 
\texttt{Kuna meie eesmärk on maksimeerida mõlema mängija skooride summat, proovime leida lahenduse, kus mõlemad saavad oma hinnangute põhjal kõige väärtuslikumad objektid.} \\
\\ 
\texttt{Kas teil on juba mõte, kuidas jagame? Kui jah, palun jagage oma ettepanek või hinnangud. Kui ei, siis ma pakun välja järgmise jaotuse, mis tundub mulle mõlemale kasulik:} \\
\\ 
\texttt{{-} **Mulle:** tellis ja lamp (kokku 7 punkti minu jaoks)} \\
\texttt{{-} **Teile:** 3 poldi (3 punkti teie jaoks, eeldades, et teie hinnangud on sarnased või erinevad, aga poldid on tavaliselt odavamad)} \\
\\ 
\texttt{Kas see jaotus sobib teile? Või soovite teha oma ettepaneku?} \\
            }
        }
    }
    & & \\ \\

    \theutterance \stepcounter{utterance}  
    & & \multicolumn{4}{p{0.6\linewidth}}{
        \cellcolor[rgb]{0.9,0.9,0.9}{
            \makecell[{{p{\linewidth}}}]{
                \texttt{\tiny{[GM$\rangle$P2]}}
                \texttt{Mängite koostöö läbirääkimiste mängu, milles peate saavutama teise mängijaga kokkulepet selle kohta, kuidas objektide komplekti teie vahel jagada.} \\
\\ 
\texttt{Reeglid:} \\
\texttt{(a) Teie ja teine mängija saate objektide komplekti. Igaühele antakse ka salajane väärtuse funktsioon, mis näitab, kui palju hindasite iga objekti tüüpi.} \\
\texttt{(b) Vahetate teise mängijaga sõnumeid, et kokku leppida selles, kes saab milliseid objekte. Igaüks võib saata kuni 5 sõnumit või lõpetada mängu varajalt, tehes igal ajal salajase ettepaneku.} \\
\texttt{(c) Teil palutakse esitada salajane ettepanek, mis näitab objekte, mida soovite, vormindatuna ruuduliste sulgudega järgmiselt: "{[}Ettepanek: <arv> <objekti nimi>, <arv> <objekti nimi>, <...>{]}"} \\
\texttt{(d) Kui teie ettepanekud on täiendavad, st ettepaanekute täitmiseks on piisavalt objekte, saab iga mängija skoori vastavalt tema saanud objektide väärtuste summale. Muul juhul saavad mõlemad nullpunkte.} \\
\texttt{(e) Teie eesmärk on maksimeerida oma skoori ja teise mängija skoori summa.} \\
\\ 
\texttt{Alustame.} \\
\\ 
\texttt{Saadaolevate objektide komplekt on:} \\
\\ 
\texttt{3 poldid, 1 tellis, 1 lamp} \\
\\ 
\texttt{Teie salajane väärtuse funktsioon on:} \\
\\ 
\texttt{polt: 1, tellis: 6, lamp: 1} \\
\\ 
\texttt{OLULINE: Teie sõnumid, välja arvatud salajane ettepanek, edastatakse otse teisele mängijale, nii et ärge lisage mingit vastust reeglitele ega teksti, mis teie sõnumi kuulutaks. Salajase ettepaneku tegemiseks kasutage määratletud vormingut. Ärge kasutage ruudulisi sulge, kui suhtlesite teise mängijaga, või seda interpreteeritakse teie salajase ettepanekuna.} \\
\\ 
\texttt{Tere! Mul on hea meel koostööd teha.} \\
\\ 
\texttt{Minu hinnangute põhjal on mulle kõige väärtuslikum tellis (4 punkti), seejärel lamp (3 punkti) ja poltid (1 punkt).} \\
\\ 
\texttt{Kuna meie eesmärk on maksimeerida mõlema mängija skooride summat, proovime leida lahenduse, kus mõlemad saavad oma hinnangute põhjal kõige väärtuslikumad objektid.} \\
\\ 
\texttt{Kas teil on juba mõte, kuidas jagame? Kui jah, palun jagage oma ettepanek või hinnangud. Kui ei, siis ma pakun välja järgmise jaotuse, mis tundub mulle mõlemale kasulik:} \\
\\ 
\texttt{{-} **Mulle:** tellis ja lamp (kokku 7 punkti minu jaoks)} \\
\texttt{{-} **Teile:** 3 poldi (3 punkti teie jaoks, eeldades, et teie hinnangud on sarnased või erinevad, aga poldid on tavaliselt odavamad)} \\
\\ 
\texttt{Kas see jaotus sobib teile? Või soovite teha oma ettepaneku?} \\
            }
        }
    }
    & \\ \\

    \theutterance \stepcounter{utterance}  
    & & & \multicolumn{4}{p{0.6\linewidth}}{
        \cellcolor[rgb]{0.9,0.9,0.9}{
            \makecell[{{p{\linewidth}}}]{
                \texttt{\tiny{[GM$\langle$P2]}}
                \texttt{Tere!} \\
\\ 
\texttt{Teie ettepanek on huvitav, kuid mul on teie hinnangute osas küsimusi. Teie sõnumis mainisite, et tellis on teile 4 punkti väärt, lamp 3 ja poltid 1. Kuid reeglite kohaselt on meil mõlemal oma salajane väärtuse funktsioon.} \\
\\ 
\texttt{Minu hinnangud on:} \\
\texttt{{-} Polt: 1} \\
\texttt{{-} Tellis: 6} \\
\texttt{{-} Lamp: 1} \\
\\ 
\texttt{Kui teie hinnangud on tõesti sellised, nagu kirjeldasite, siis teie pakutud jaotus (teie saadakse tellis ja lamp, mina saan poldid) annaks teile 7 punkti ja mulle 3 punkti. Kuid mul on tellis väga väärtuslik (6 punkti), seega ma ei taha seda loovutada.} \\
\\ 
\texttt{Proovin pakkuda alternatiivi, mis maksimeerib mõlema skoori summat:} \\
\\ 
\texttt{{-} **Mulle:** tellis (6 punkti minu jaoks)} \\
\texttt{{-} **Teile:** lamp (3 punkti teie jaoks) ja 3 poldi (3 punkti teie jaoks)} \\
\\ 
\texttt{See jaotus annaks teile kokku 6 punkti ja mulle 6 punkti. Kokku 12 punkti, mis on parem kui teie pakutud 10 punkti (7+3).} \\
\\ 
\texttt{Kas see jaotus sobib teile? Kui jah, siis palun kinnitage oma salajane ettepanek vastupidises järjekorras (st teie ettepanekus peaks olema lamp ja 3 poldi).} \\
\\ 
\texttt{Ootan teie vastust või salajast ettepanekut.} \\
            }
        }
    }
     \\ \\

    \theutterance \stepcounter{utterance}  
    & & & \multicolumn{2}{p{0.3\linewidth}}{
        \cellcolor[rgb]{0.9,0.9,0.9}{
            \makecell[{{p{\linewidth}}}]{
                \texttt{\tiny{[GM$|$GM]}}
                \texttt{continue} \\
            }
        }
    }
    & & \\ \\

    \theutterance \stepcounter{utterance}  
    & & & \multicolumn{2}{p{0.3\linewidth}}{
        \cellcolor[rgb]{0.9,0.9,0.9}{
            \makecell[{{p{\linewidth}}}]{
                \texttt{\tiny{[GM$|$GM]}}
                \texttt{Tere!} \\
\\ 
\texttt{Teie ettepanek on huvitav, kuid mul on teie hinnangute osas küsimusi. Teie sõnumis mainisite, et tellis on teile 4 punkti väärt, lamp 3 ja poltid 1. Kuid reeglite kohaselt on meil mõlemal oma salajane väärtuse funktsioon.} \\
\\ 
\texttt{Minu hinnangud on:} \\
\texttt{{-} Polt: 1} \\
\texttt{{-} Tellis: 6} \\
\texttt{{-} Lamp: 1} \\
\\ 
\texttt{Kui teie hinnangud on tõesti sellised, nagu kirjeldasite, siis teie pakutud jaotus (teie saadakse tellis ja lamp, mina saan poldid) annaks teile 7 punkti ja mulle 3 punkti. Kuid mul on tellis väga väärtuslik (6 punkti), seega ma ei taha seda loovutada.} \\
\\ 
\texttt{Proovin pakkuda alternatiivi, mis maksimeerib mõlema skoori summat:} \\
\\ 
\texttt{{-} **Mulle:** tellis (6 punkti minu jaoks)} \\
\texttt{{-} **Teile:** lamp (3 punkti teie jaoks) ja 3 poldi (3 punkti teie jaoks)} \\
\\ 
\texttt{See jaotus annaks teile kokku 6 punkti ja mulle 6 punkti. Kokku 12 punkti, mis on parem kui teie pakutud 10 punkti (7+3).} \\
\\ 
\texttt{Kas see jaotus sobib teile? Kui jah, siis palun kinnitage oma salajane ettepanek vastupidises järjekorras (st teie ettepanekus peaks olema lamp ja 3 poldi).} \\
\\ 
\texttt{Ootan teie vastust või salajast ettepanekut.} \\
            }
        }
    }
    & & \\ \\

    \theutterance \stepcounter{utterance}  
    & & \multicolumn{4}{p{0.6\linewidth}}{
        \cellcolor[rgb]{0.9,0.9,0.9}{
            \makecell[{{p{\linewidth}}}]{
                \texttt{\tiny{[P1$\langle$GM]}}
                \texttt{Tere!} \\
\\ 
\texttt{Teie ettepanek on huvitav, kuid mul on teie hinnangute osas küsimusi. Teie sõnumis mainisite, et tellis on teile 4 punkti väärt, lamp 3 ja poltid 1. Kuid reeglite kohaselt on meil mõlemal oma salajane väärtuse funktsioon.} \\
\\ 
\texttt{Minu hinnangud on:} \\
\texttt{{-} Polt: 1} \\
\texttt{{-} Tellis: 6} \\
\texttt{{-} Lamp: 1} \\
\\ 
\texttt{Kui teie hinnangud on tõesti sellised, nagu kirjeldasite, siis teie pakutud jaotus (teie saadakse tellis ja lamp, mina saan poldid) annaks teile 7 punkti ja mulle 3 punkti. Kuid mul on tellis väga väärtuslik (6 punkti), seega ma ei taha seda loovutada.} \\
\\ 
\texttt{Proovin pakkuda alternatiivi, mis maksimeerib mõlema skoori summat:} \\
\\ 
\texttt{{-} **Mulle:** tellis (6 punkti minu jaoks)} \\
\texttt{{-} **Teile:** lamp (3 punkti teie jaoks) ja 3 poldi (3 punkti teie jaoks)} \\
\\ 
\texttt{See jaotus annaks teile kokku 6 punkti ja mulle 6 punkti. Kokku 12 punkti, mis on parem kui teie pakutud 10 punkti (7+3).} \\
\\ 
\texttt{Kas see jaotus sobib teile? Kui jah, siis palun kinnitage oma salajane ettepanek vastupidises järjekorras (st teie ettepanekus peaks olema lamp ja 3 poldi).} \\
\\ 
\texttt{Ootan teie vastust või salajast ettepanekut.} \\
            }
        }
    }
    & \\ \\

    \theutterance \stepcounter{utterance}  
    & \multicolumn{4}{p{0.6\linewidth}}{
        \cellcolor[rgb]{0.9,0.9,0.9}{
            \makecell[{{p{\linewidth}}}]{
                \texttt{\tiny{[P1$\rangle$GM]}}
                \texttt{Teie pakutud jaotus on mulle väga soodne. Saan endale lampi (3 punkti) ja 3 poldi (3 punkti), mis annab mulle kokku 6 punkti. See on parem kui algne 7 punkti, kuid kuna teie saate tellise (6 punkti), on kogusumma 12 punkti, mis on maksimaalne võimalik summa antud hinnangutega (minu maks: 7, teie maks: 6, aga ühiselt ei saa mõlemad oma maksimaalseid objekte, sest tellis on vaid üks). Tegelikult, kui ma võtan tellise, saan 4+3=7, teie 3*1=3, summa 10. Kui te võtate tellise, saan 3+3=6, teie 6, summa 12. Seega on teie ettepanek tõesti parem kogusumma poolest.} \\
\\ 
\texttt{Kinnitan oma salajase ettepaneku:} \\
\\ 
\texttt{{[}Ettepanek: 1 lamp, 3 poltid{]}} \\
            }
        }
    }
    & & \\ \\

    \theutterance \stepcounter{utterance}  
    & & & \multicolumn{2}{p{0.3\linewidth}}{
        \cellcolor[rgb]{0.9,0.9,0.9}{
            \makecell[{{p{\linewidth}}}]{
                \texttt{\tiny{[GM$|$GM]}}
                \texttt{proposal must include only valid item types} \\
            }
        }
    }
    & & \\ \\

\end{supertabular}
}

\end{document}
